\section{Discussions}

We have introduced Anchored E-Watermarking, a novel framework that bridges the gap between optimal sampling and anytime-valid inference in statistical watermarking. By shifting the detection paradigm from p-values to e-values, we addressed the critical limitation of fixed-horizon testing, enabling valid optional stopping without compromising Type-I error guarantees.

Moreover, we characterized the optimal e-value with respect to the worst-case log-growth rate and derived the optimal expected stopping time, providing a rigorous foundation for watermarking in the presence of an anchor distribution. Empirically, our results on real-world language models demonstrate that this principled approach translates into substantial gains in efficiency. Our method identifies watermarked content with significantly fewer tokens than state-of-the-art heuristics while preserving generation quality.

As the first application of e-values to statistical watermarking, this framework opens new avenues for efficient detection mechanisms, with future works including extension to more flexible anchor distributions or investigating the game-theoretic implications of e-watermarking against incentivized adversaries.

